% ---
% Reemplaza líneas 1803-1844 de Docs/Technical_Report.tex
% Decisiones: D-016, D-033
% ---

\subsubsection{Índice Davies-Bouldin}

El índice Davies-Bouldin (DB) \cite{davies1979cluster}
mide la calidad de la separación entre clusters en un
espacio de representación, combinando la dispersión
intra-clase con la separación inter-clase. Para el
problema de clasificación binaria (benigno / maligno),
las dos clases diagnósticas se definen como clusters
predeterminados ($C = 2$). El índice se define como:

\begin{equation}
  \text{DB} =
  \frac{1}{C}\sum_{i=1}^{C}
  \max_{j \neq i}
  \left\{
  \frac{s_i + s_j}{d(c_i, c_j)}
  \right\}
  \label{eq:davies_bouldin}
\end{equation}

\noindent donde $C = 2$ es el número de clases,
$s_i$ es la dispersión promedio de los puntos de la
clase $i$ respecto a su centroide $c_i$:

\begin{equation}
  s_i = \frac{1}{|C_i|}
  \sum_{\mathbf{z} \in C_i}
  \|\mathbf{z} - c_i\|_2
  \label{eq:dispersion}
\end{equation}

\noindent y $d(c_i, c_j) = \|c_i - c_j\|_2$ es la
distancia euclidiana entre los centroides de las clases
$i$ y $j$. Valores bajos del índice DB indican clusters
compactos y bien separados; el valor óptimo es 0.
Esta métrica es independiente del clasificador y puede
aplicarse directamente sobre el vector de embeddings
$\mathbf{e}(\mathbf{x}, \boldsymbol{\theta})$ de las
condiciones cuánticas o sobre los vectores transformados
de las condiciones clásicas. Al depender de distancias
euclidianas invariantes ante rotaciones ortogonales,
el índice DB genera el mismo valor para C1 y C2.
